
\documentclass[11pt]{article}
\usepackage[margin=1in]{geometry}
\usepackage{amsmath,amssymb,booktabs,hyperref,graphicx,parskip,enumitem}
\usepackage{xcolor}
\hypersetup{colorlinks=true,linkcolor=blue!50!black,urlcolor=blue!50!black}
\setlist{nosep}

\title{Replication Report: FDTD: Solving 1+1D Delay PDE in Parallel\\[2pt]\large OSTI 1475143 \textbar{} Coverage 8/10, Agreement 8/10}
\author{Rick Stevens \and Ollie (AI Assistant)\\\small Argonne National Laboratory}
\date{2026-04-27}

\begin{document}\maketitle

\section{Source paper}
\begin{itemize}
\item \textbf{Title:} FDTD: Solving 1+1D Delay PDE in Parallel
\item \textbf{Authors:} Zheng, Goldschmidt, Fan
\item \textbf{Year:} 2018
\item \textbf{Domain:} Quantum Optics / Numerical PDE
\item \textbf{Identifier:} OSTI 1475143
\end{itemize}


\section{Replication objective}
The central claim of the paper concerns quantum optics / numerical pde. Our objective was to reproduce the headline numerical/qualitative results within the constraints of the AI-driven Replication Project (24--72\,h, commodity GPU/CPU compute, no author contact). A replication plan is archived at \texttt{1475143-FDTD-solving-1+1D-delay-PDE-in/replication\_plan.pdf}.

\section{Methodology}
We followed the standard Replication Project workflow: ingest the source PDF, extract method and data pointers, draft a replication plan, attempt the smallest end-to-end pipeline that produces a comparable number, then scale up where compute and data permit. Tools used in this replication are catalogued in the meta-paper's computational tools inventory (see \texttt{COMPUTATIONAL\_TOOLS\_INVENTORY.md}).

This paper went through the Tier-Lift pipeline; supplementary notes at \texttt{.openclaw/workspace/24h-progress/batch1-pde/1475143.md}.

\subsection*{Paper-directory README excerpt}
\small
\par\medskip\textbf{FDTD: solving 1+1D delay PDE in parallel}\par
\par$\bullet$\ **OSTI ID:** 1475143
\par$\bullet$\ **Rank:** \#6
\par$\bullet$\ **Replication Score:** 10/10

\par\medskip\textbf{\# Why This Paper}\par
This paper is a perfect candidate as it describes a fully computational workflow (FDTD for delay PDEs), uses structured simulation data, and provides explicit algorithmic and benchmarking details for quantitative validation.

\par\medskip\textbf{\# Replication Plan}\par
AI would implement the FDTD algorithm for the specified delay PDE, generate simulation data for various parameter sets, and benchmark parallelization approaches, validating results against quantitative metrics and example outputs.

\par\medskip\textbf{\# Status}\par
\par$\bullet$\ [ ] Paper reviewed
\par$\bullet$\ [ ] Data identified
\par$\bullet$\ [ ] Code implemented
\par$\bullet$\ [ ] Results validated
\normalsize

The replication was driven by an OpenClaw agent loop (Claude Opus / GPT-5 class models) issuing shell, Python, and (where applicable) GPU jobs. Compute usage and wall-time for individual papers are summarised in the meta-paper's resource table; not separately recorded for this paper unless noted in the Tier-Lift artefact. Sub-agents were spawned for replication-plan generation and (for papers in batches 1--4) for tier-lift extension runs.

\section{What was reproduced}
\begin{itemize}
\item Numba-JIT square-rule FDTD (52x speedup)
\item Bit-exact Fig 5 max|psi| at t/Delta=99/500/999/39999
\item 4-worker parameter-sweep parallelism
\item Fig 6 antibunching (g2=0.001) + bunching (g2=2.13)
\item Fig 7 antibunching (g2=0.007) + oscillations
\end{itemize}

\section{What was missing or incomplete}
\begin{itemize}
\item Fig 8 non-Markovian (125 GB memory wall)
\item Ref [8] scattering-theory g2 comparison
\item True intra-timestep parallelism
\item GPU/CUDA port
\end{itemize}
The dominant blocker categories for this paper, mapped onto the meta-paper's 9-category friction taxonomy (\S4), are listed in \S\ref{sec:friction} below.

\section{Quantitative agreement}
Per-quantity numerical comparisons are not separately tabulated in the structured evaluation record for this paper beyond the agreement rationale below; where the rationale references specific numbers, they are reproduced verbatim. A full numerical comparison table is recorded in the per-replication artefacts (see \S\ref{sec:repro}). For papers below 8/8, the agreement rationale identifies the specific quantities that diverge.

\paragraph{Agreement rationale.} Fig 5 max|psi| matches reference to 6+ digits at all snapshots. Fig 6: g2(0)=0.0007 at resonance (antibunching), 2.13 off-resonance (bunching), matching paper within \textasciitilde{}5\%. Fig 7: g2(0)=0.007 at resonance, 0.539 off-resonance (\textasciitilde{}10\% of paper visual). Numba throughput 7800 steps/s same order as paper's parallel implementation. All physics trends correct.

\section{Score and friction-point tagging}\label{sec:friction}
\begin{itemize}
\item \textbf{Coverage:} 8/10 --- Adds Numba-JIT FDTD core (52x speedup, 7808 steps/s, bit-identical to Python reference at all Fig 5 snapshots) and multi-grid parameter-sweep parallel benchmark (1.31x at 4 workers). Now reproduces 5/6 paper figures (Fig 5/6/7 + parallel-performance theme). Fig 8 still memory-bound (\textasciitilde{}125 GB). Prior Python solver + analytical BCs + g2 correlation from Figs 6-7 retained.
\item \textbf{Agreement:} 8/10 --- Fig 5 max|psi| matches reference to 6+ digits at all snapshots. Fig 6: g2(0)=0.0007 at resonance (antibunching), 2.13 off-resonance (bunching), matching paper within \textasciitilde{}5\%. Fig 7: g2(0)=0.007 at resonance, 0.539 off-resonance (\textasciitilde{}10\% of paper visual). Numba throughput 7800 steps/s same order as paper's parallel implementation. All physics trends correct.
\end{itemize}

\noindent\textbf{Friction points encountered (taxonomy tags):}
\begin{itemize}
\item None recorded.
\end{itemize}

\section{Follow-on questions}
\begin{itemize}
\item Q1: Port to CUDA -- does antidiagonal sweep beat sequential at GPU clock?
\item Q2: Scattering-theory g2 reference for O(Delta\textasciicircum{}2) convergence validation?
\item Q3: Out-of-core (zarr/dask) for Fig 8 regime?
\item Q4: Phase diagram g2(0) vs (k0a, omega0/Gamma)?
\item Q5: Finite-reflectivity mirror r\textless{}1 -- Markov transition threshold?
\end{itemize}

\section{Reproducibility pointers}\label{sec:repro}
\begin{itemize}
\item \textbf{Paper directory:} \texttt{1475143-FDTD-solving-1+1D-delay-PDE-in}
\item \textbf{Replication artefacts:}
\begin{itemize}
\item Replication artifacts in paper directory.
\end{itemize}
\item \textbf{Replication-dir hint from JSONL:} \texttt{\textasciitilde{}/projects/replicate-fdtd/}
\item \textbf{Status:} tier\_lift\_2026\_04\_26
\end{itemize}

To re-run, change directory into the paper directory above and consult its \texttt{README.md} for the entry-point script. Most replications follow the convention \texttt{replication/run\_*.sh} or \texttt{replication/<subdir>/run.py}.

\end{document}
